\chapter{Experiments and Results}
\label{chap:experiments}

\section{Overview}
This chapter reports the outcomes of the experiments declared in profile \texttt{\ThesisProfileId}.
Every value below is a macro generated from the signed result bundle; tables and figures are
produced by the same pipeline.

\input{generated/auto_tables.tex}

\ifAutoHasComparison
\section{Comparison with baseline methods}
Table~\ref{tab:auto-comparison} summarises all metrics. On \MetricPrimaryName, the proposed
method (\MetricPrimaryProposedName) achieves \MetricPrimaryProposedMean{} $\pm$
\MetricPrimaryProposedStd{} \MetricPrimaryUnit{} (confidence interval \MetricPrimaryProposedCiLow{} to
\MetricPrimaryProposedCiHigh), whereas the strongest baseline (\MetricPrimaryRunnerUpName) reaches
\MetricPrimaryRunnerUpMean{} $\pm$ \MetricPrimaryRunnerUpStd. The difference of \DiffPrimaryAbs{}
(\DiffPrimaryRelPct\,\% relative) is assessed with the \StatPrimaryTest{} ($p$ \StatPrimaryPval;
\StatPrimaryEffectName{} \StatPrimaryEffect; \StatPrimaryN{} runs per method). Across all methods,
the \StatPrimaryOmnibusTest{} yields $p$ \StatPrimaryOmnibusPval.
\fi

\ifAutoHasAblation
\section{Ablation study}
Table~\ref{tab:auto-ablation} reports the ablation configurations. The full system reaches
\MetricPrimaryProposedMean{} $\pm$ \MetricPrimaryProposedStd{} \MetricPrimaryUnit{} on
\MetricPrimaryName, against \MetricPrimaryRunnerUpMean{} for the best reduced configuration
(\MetricPrimaryRunnerUpName), a gain of \DiffPrimaryAbs{} (\DiffPrimaryRelPct\,\%; \StatPrimaryTest,
$p$ \StatPrimaryPval). The omnibus \StatPrimaryOmnibusTest{} across configurations gives
$p$ \StatPrimaryOmnibusPval.
\fi

\ifAutoHasRsm
\section{Response surface optimisation}
The design runs are listed in Table~\ref{tab:auto-design}. The fitted model
(Table~\ref{tab:auto-coef}) reads
\begin{equation}
  \StatRsmPrimaryEquation
\end{equation}
with \StatRsmPrimaryFactorLegend. The ANOVA in Table~\ref{tab:auto-anova} gives
$R^{2}$ = \StatRsmPrimaryRsq{} and adjusted $R^{2}$ = \StatRsmPrimaryRsqAdj{}
($F$ = \StatRsmPrimaryFvalue, $p$ \StatRsmPrimaryPval), and the lack-of-fit test returns
$p$ \StatRsmPrimaryLackOfFitPval. The stationary point is a \StatRsmPrimaryOptNature{} located at
\StatRsmPrimaryOptPoint, with a predicted response of \StatRsmPrimaryOptResponse{}
\MetricPrimaryUnit. Figure~\ref{fig:auto-rsm} visualises the surface.
\fi

\section{Figures}
\input{generated/auto_figures.tex}
